\documentclass[12pt]{article}
\usepackage[utf8]{inputenc}
\usepackage[T1]{fontenc}
\usepackage[margin=1in]{geometry}
\usepackage{setspace}
\usepackage{xcolor}
\usepackage[colorlinks=true, urlcolor=blue!60!black, linkcolor=blue!60!black]{hyperref}
\setlength{\parindent}{0em}
\setlength{\parskip}{0.6em}
\renewcommand{\baselinestretch}{1.0}
\begin{document}
\begin{center}
{\LARGE \textit{Abbott Laboratories v. Gardner}, 387 U.S. 136 (1967)\par}
\vspace{0.5em}
{\large Case Brief --- Two-Page Edition\par}
\end{center}

\section{Memory Jogger}
Final agency action is presumptively reviewable, and a pre-enforcement challenge is ripe when the issue is fit for decision and delay would cause hardship (\href{https://supreme.justia.com/cases/federal/us/387/136/}{Justia}).

\section{Detailed Case Facts}
An FDA rule required prescription drug makers to print a drug's generic ``established name'' with every appearance of its trade name---the ``every time'' rule. Violations exposed makers to seizures, injunctions, and prosecution.

\section{Procedural History}
Before enforcement, 37 manufacturers sued in the District of Delaware, which held the rule unauthorized. The Third Circuit reversed, finding review precluded and no case or controversy.

\section{Judicial Votes}
Reversed and remanded, 5--3. Harlan for the Court (Black, Douglas, Stewart, White); Fortas, with Warren and Clark, and Clark separately, dissenting; Brennan not participating.

\section{Holding}
Review of final agency action is presumed absent ``clear and convincing evidence'' of contrary congressional intent. The challenge was ripe: the issue was purely legal, the rule was final, and withholding review would impose immediate hardship.

\section{Analysis of Opinions}
\textbf{Harlan:} The APA embodies a ``basic presumption of judicial review,'' and the Act's special review provision did not bar other remedies. Ripeness weighs fitness and hardship; the formal, immediately effective rule forced makers to choose between costly relabeling and risking prosecution. \textbf{Fortas:} the Court gave district judges a ``roving commission to halt the regulatory process''; challenges belong in enforcement proceedings.

\section{Examples: Future Applications}
\textbf{Same-Side:} A final rule forcing data brokers to delete location data daily, with heavy penalties, is ripe. \textbf{Opposite-Side:} A rule saying an agency ``may'' audit brokers in its discretion is unripe (\textit{Toilet Goods}). \textbf{Fence-Sitter:} Guidance backed by warning letters creates pressure but may not be final agency action.

\section{Critique}
\textit{Abbott} blurred Article III and prudential ripeness, a gap exposed after \textit{Lexmark}. Critics echoing Fortas say it shifts policymaking from agencies to judges.

\section{Key Quotations}
\begin{quote}
``Without undertaking to survey the intricacies of the ripeness doctrine it is fair to say that its basic rationale is to prevent the courts, through avoidance of premature adjudication, from entangling themselves in abstract disagreements over administrative policies, and also to protect the agencies from judicial interference until an administrative decision has been formalized and its effects felt in a concrete way by the challenging parties. The problem is best seen in a twofold aspect, requiring us to evaluate both the fitness of the issues for judicial decision and the hardship to the parties of withholding court consideration.''
\end{quote}
\noindent --- Justice Harlan, Opinion of the Court, 387 U.S. at 148--149.

\end{document}
